% TikZ 图：黑体辐射与量子化思想来源
% Day 27：近代物理与原子物理入门
\begin{tikzpicture}[>=Latex,scale=0.92,every node/.style={font=\scriptsize}]
  \draw[->,thick] (0,0) -- (7.1,0) node[right] {波长 $\lambda$};
  \draw[->,thick] (0,0) -- (0,3.95) node[above] {光谱辐射强度};

  \draw[thick,blue!70!black,smooth]
    plot coordinates {(0.25,0.04) (0.7,0.38) (1.25,1.20) (1.95,2.05) (2.75,1.55) (3.65,0.90) (4.75,0.43) (6.2,0.14)};
  \draw[thick,red!80!black,smooth]
    plot coordinates {(0.22,0.08) (0.55,0.75) (0.95,2.25) (1.45,3.25) (2.05,2.45) (2.9,1.35) (4.1,0.55) (5.8,0.18)};

  \draw[dashed,blue!70!black] (1.95,0) -- (1.95,2.05);
  \draw[dashed,red!80!black] (1.45,0) -- (1.45,3.25);
  \node[blue!70!black,align=center] at (3.85,2.05) {温度较低\\峰值波长较长};
  \node[red!80!black,align=center] at (4.85,3.15) {温度较高\\峰值向短波移动};

  \draw[->,orange!85!black,thick] (2.05,3.45) -- (1.50,3.45);
  \node[orange!85!black,align=center] at (2.95,3.45) {升温};

  \draw[gray!70,dashed,thick]
    (3.60,3.75) .. controls (3.95,2.55) and (4.50,0.65) .. (5.85,0.18);
  \node[gray!70!black,align=center] at (6.15,1.55) {经典理论\\短波端发散};

  \node[draw,rounded corners=2pt,fill=yellow!12,align=center,inner sep=4pt] at (4.55,0.75)
    {普朗克：能量交换\\按 $E=h\nu$ 一份一份进行};
\end{tikzpicture}
